\chapter{Casi di Studio: Dataset DEAP}

\section{Introduzione a DEAP}
Il dataset \textbf{DEAP} (Database for Emotion Analysis using Physiological Signals) è uno standard di fatto per l'Affective Computing basato su segnali fisiologici e video facciali. Creato per sviluppare sistemi di riconoscimento delle emozioni, ha introdotto concetti essenziali riguardo all'acquisizione multimodale.

Le caratteristiche principali del protocollo di induzione emotiva usato in DEAP:
\begin{itemize}
    \item \textbf{Stimoli usati:} Video musicali da 1 minuto scelti per indurre emozioni uniformemente distribuite nel piano di Russell.
    \item \textbf{Segnali Acquisiti:}
    \begin{itemize}
        \item Segnali fisiologici classici periferici: GSR, ECG, Respiro, Temperatura, EOG, EMG.
        \item Segnali Cerebrali: \textbf{EEG a 32 canali}.
        \item Segnali Visivi: \textbf{Face Video} dei partecipanti.
    \end{itemize}
\end{itemize}

\section{Formato di Valutazione (Self-Assessment)}
Subito dopo ogni video, il soggetto valutava il proprio stato emotivo usando il sistema visivo dei manichini \textbf{SAM (Self-Assessment Manikins)}, esprimendosi sulle seguenti dimensioni (su scala continua 1-9):
\begin{itemize}
    \item \textbf{Arousal (Attivazione):} da Calmo a Eccitato.
    \item \textbf{Valence (Valenza):} da Triste/Negativo a Felice/Positivo.
    \item \textbf{Dominance (Controllo):} da Sottomesso a Dominante.
    \item \textbf{Liking (Gradimento):} quanto al soggetto è piaciuto il video.
\end{itemize}

\section{Modality Fusion}
Una delle conclusioni più importanti emerse dallo studio sul paper DEAP riguarda la \textbf{Modality Fusion} (fusione multimodale). Unendo \textbf{feature estratte dall'EEG} assieme a \textbf{feature estratte dai segnali periferici e dalla MCA (Multimedia Content Analysis degli stimoli video)}, le prestazioni dei classificatori sono notevolmente superiori rispetto all'uso di un singolo biosegnale.

